\section{Steinberg groups and faithful strips}\label{sec:steinberg}

We next introduce the group that records elementary matrix operations
before imposing any relations beyond the Steinberg relations. Its kernel
over matrices is the part that finite degree truncations must detect.
All algebras in this section are unital associative $\F_2$-algebras.

For such an algebra $B$ and $n\ge3$, the group $\St_n(B)$ has generators
$x_{ij}(a)$, where $a\in B$ and $1\le i,j\le n$ with $i\ne j$, and
relations
\begin{align}
 x_{ij}(a)x_{ij}(b)&=x_{ij}(a+b),\label{eq:root-addition}\\
 [x_{ij}(a),x_{lt}(b)]&=1
       &&(j\ne l,\ i\ne t),\label{eq:root-commuting}\\
 [x_{ij}(a),x_{jt}(b)]&=x_{it}(ab)
       &&(i,j,t\text{ distinct}).\label{eq:root-product}
\end{align}
Our convention is $[X,Y]=XYX^{-1}Y^{-1}$.
In particular, $x_{ij}(0)=1$ and each root element has square $1$.
The elementary matrix $e_{ij}(a)$ is the identity with $a$ added in
position $(i,j)$. These matrices satisfy
\eqref{eq:root-addition}--\eqref{eq:root-product}, giving an epimorphism
\[
 \pi_n:\St_n(B)\longrightarrow\E_n(B)
       :=\langle e_{ij}(a):a\in B,\ i\ne j\rangle.
\]
Write $J_n(B)=\ker\pi_n$. We use the same presentation with a countable
index set to define $\St(B)$, and put
\[
 \E(B)=\varinjlim_n\E_n(B),\qquad
 J(B)=\ker\bigl(\St(B)\longrightarrow\E(B)\bigr).
\]
Thus $\St(B)=\varinjlim_n\St_n(B)$: each word and each finite derivation
uses only finitely many indices. This is a direct limit of discrete
groups. No completion is involved. The fixed-rank injectivity needed for
$R$ will be proved separately in Section~\ref{sec:stability}.

The presentation and its relation to algebraic $K$-theory are classical;
see \cite{Milnor,WeibelK}. We will use the following elementary
consequences repeatedly. They also explain when an identity of matrices
does imply the corresponding identity in the Steinberg group.

\begin{lemma}[Faithful strips]\label{lem:strip}
Let $S,T$ be disjoint finite sets of coordinates. The subgroup generated
by $x_{ij}(a)$ with $i\in S$, $j\in T$ maps isomorphically onto the
additive group of matrices supported in $S\times T$ under $u\mapsto
\pi(u)-\Id$. A word supported on $S$ normalizes this subgroup and acts
on its entries by left multiplication by its matrix on $S$.
On the reverse strip $T\times S$, its action is right multiplication
by the inverse matrix on $S$.
\end{lemma}

\begin{proof}
All roots in $S\times T$ commute by \eqref{eq:root-commuting}.
Root addition collects a word into one factor per coordinate pair.
The corresponding entry matrices have pairwise zero products, so the
matrix image reads each collected coefficient directly. This proves
injectivity and identifies the group law with addition.

For one root supported on $S$, conjugation on an outgoing strip follows
from \eqref{eq:root-product}: it performs the corresponding elementary
row addition. The same relations on the reverse strip perform the
inverse elementary column operation. Iterating over the factors of a
word gives the stated left and right actions, with their prescribed
order over the possibly noncommutative algebra $B$.
\end{proof}

\begin{lemma}[Stable centrality]\label{lem:stable-center}
For every $B$, the group $J(B)$ is central in $\St(B)$.
\end{lemma}

\begin{proof}
Let $z\in J(B)$ and choose a finite set $S$ containing the support of a
representative of $z$. To show that $z$ commutes with a specified root,
enlarge $S$ to contain both indices of that root. Choose $h\notin S$.
By Lemma~\ref{lem:strip}, $z$ acts trivially on the strips $S\times\{h\}$
and $\{h\}\times S$, since its matrix image is the identity. For
$i,j\in S$ with $i\ne j$, the identity
\[
 x_{ij}(a)=[x_{ih}(a),x_{hj}(1)]
\]
therefore implies that $z$ commutes with $x_{ij}(a)$. These roots
generate the group.
\end{proof}

\begin{lemma}[Coordinate permutations]\label{lem:permutations}
In characteristic two, every permutation of a finite set of coordinates
is induced on roots by conjugation in the Steinberg group after adjoining
an outside coordinate if necessary. In particular, relabeling the
finitely many coordinates in a representative of $z\in J(B)$ leaves
the element $z$ unchanged in $\St(B)$.
\end{lemma}

\begin{proof}
For a transposition of coordinates $a,b$, use
\[
 p=x_{ab}(1)x_{ba}(1)x_{ab}(1).
\]
On a column strip from $\{a,b\}$ to an outside coordinate, its three
elementary row additions exchange the two entries over $\F_2$.
Lemma~\ref{lem:strip} gives this as an exact group action. The reverse
strip is exchanged in the same way. Roots with both indices outside
$\{a,b\}$ commute with $p$. The remaining two root positions can be
written, with $c\notin\{a,b\}$, as
\[
 x_{ab}(u)=[x_{ac}(u),x_{cb}(1)],\qquad
 x_{ba}(u)=[x_{bc}(u),x_{ca}(1)].
\]
The strip actions exchange these as well. Products of transpositions
give arbitrary finite permutations. Any injective relabeling of a
finite set extends to a permutation of a larger finite set, so the
last assertion follows from Lemma~\ref{lem:stable-center}.
\end{proof}

This proof of Lemma~\ref{lem:permutations} is the argument of
\cite[Lemma~6.3]{Companion}; its exact action on roots, rather than only
on their matrix images, will matter below.

\begin{lemma}[A commutator identity]\label{lem:commutator}
In any group, let $A=[X,Y]$ and $B'=[Y,Z]$. If
\[
 [X,Z]=[A,X]=[B',Z]=[A,B']=1,
\]
then $[A,Z]=[X,B']$. This common commutator commutes with
$A,Z,X,B'$.
\end{lemma}

\begin{proof}
The subgroups $H=\langle A,Z\rangle$ and $K=\langle X,B'\rangle$
commute elementwise by the four hypotheses. Conjugating $[X,Z]=1$
by $Y$ gives
\[
 [A^{-1}X,B'Z]=1.
\]
Separating the commuting $H$ and $K$ factors yields
$[A^{-1},Z][X,B']=1$. Thus $[X,B']=[A^{-1},Z]^{-1}$ lies in
$H\cap K$, which is central in both $H$ and $K$. Centrality in $H$
then gives $[A,Z]=[A^{-1},Z]^{-1}=[X,B']$.
\end{proof}

Finally, $\St_n(B)$ and $\E_n(B)$ are perfect for $n\ge3$:
each generating root is the commutator in \eqref{eq:root-product}
with one coefficient equal to $1$. If $B$ is infinite, both groups are
infinite, since distinct coefficients at a fixed root position give
distinct elementary matrices.
